\chapter{Semantic Changes and Warnings}
\mpitermtitleindex{semantic changes}
\label{chap:semantic-changes}

This chapter lists semantic changes that have been introduced into the \MPI/ standard as well as
warnings that could potentially impact program behavior. In addition to those listed here,
Chapter~\ref{chap:removed} also lists changes and backward incompatibilities caused by removing
interfaces. Unlike Chapter~\ref{chap:removed}, the changes in this chapter did not go through a deprecation
process.

\section{Semantic Changes}
\label{sec:semantic-changes:changes}

This section describes semantics that have changed in a way that would potentially cause an \MPI/
program to behave differently when using this version of the \MPI/ standard without changing the program's code.

\subsection{Semantic Changes Starting in \texorpdfstring{\mpiivdoto/}{MPI-4.0}}
\label{subsec:semantic-changes:changes:since40}

\begin{itemize}
\item
\mpifunc{MPI\_COMM\_DUP} and \mpifunc{MPI\_COMM\_IDUP} no longer propagate info hints from the input
communicator to the output communicator. This behavior can be achieved using
\mpifunc{MPI\_COMM\_DUP\_WITH\_INFO} and \mpifunc{MPI\_COMM\_IDUP\_WITH\_INFO}.


\deprecatedsep
\item
The default communicator where errors are raised when not involving a
communicator, window, or file was changed from \mpiconst{MPI\_COMM\_WORLD} to
\mpiconst{MPI\_COMM\_SELF}.

\end{itemize}

\section{Additional Warnings}
\label{sec:semantic-changes:warnings}

This section describes additional changes that could potentially cause a program that relies on the semantics described in a previous version of the \MPI/ standard to behave differently than with this version of \MPI/. The changes in this section are limited in scope and unlikely to impact most programs.

\subsection{Warnings Starting in \texorpdfstring{\mpiivdoti/}{MPI-4.1}}
\label{subsec:semantic-changes:warnings:since41}
\label{sec:incompatible:since41} % to be substitued as soon as both the change-log entry and this issue is merged.

Implementations are no longer allowed to implement \mpifunc{MPI\_WTICK}, \mpifuncskip{PMPI\_WTICK}, \mpifunc{MPI\_WTIME}, and \mpifuncskip{PMPI\_WTIME}
as well as handle conversion functions as macros (Sections~\ref{sec:timers}
and~\ref{sec:misc-handleconvert}).
This should not impact applications but may require changes in some implementations.

\subsection{Warnings Starting in \texorpdfstring{\mpiivdoto/}{MPI-4.0}}
\label{subsec:semantic-changes:warnings:since40}
\label{sec:incompatible:since40} % to be substitued as soon as both the change-log entry and this issue is merged.

The limit for length of \MPI/ identifiers was removed.
Prior to \MPIIVDOTO/, \MPI/ identifiers were limited to 30 characters
(31 with the profiling interface).
This limitation was initially introduced to avoid exceeding the limit on some compilation systems.

\begin{rationale}
For Fortran, this limit was already relaxed for the Fortran specific function names,
see \sectionref{sec:f90:linker-names}, and the Fortran language specification 2003
requires support for a minimum of 63 characters for internal and external identifiers.
Starting with the ISO/IEC 9899:1999 C programming language standard,
support for a minimum of 63 characters is required for internal identifiers,
but only 31 characters are required to be significant for external identifiers.
At the time of the release of \MPIIVDOTO/, most or nearly all compilers allow external
identifiers longer than 31 characters. Therefore, the restriction is removed.
\end{rationale}

\begin{users}
This affects users only if they store \MPI/ identifiers into fixed sized strings.
\end{users}
